\documentclass[aps,prd,preprint,superscriptaddress,amsmath,amssymb]{revtex4-2}
\emergencystretch=2.5em
\usepackage{graphicx}
\usepackage{bm}
\usepackage{dcolumn}
\usepackage{amsthm}
\usepackage{booktabs}
\usepackage{color}

\newtheorem{postulate}{Postulate}[section]
\newtheorem{theorem}{Theorem}[section]

\begin{document}

\title{The Ontology of the Universe in the Vibrating Coordinate Frame}

\author{Yuan-Jie Liu}
\email{yuanjieliu65@gmail.com}

\date{\today}

\begin{abstract}
This paper proposes a new ontology of the universe: the cosmic
substrate is a vast collective of vibration frequencies and
dimensions, and physics is the description process of that field in
the vibrating coordinate frame.  In this ontology, time, order, band
number, dimension, width, and the ``instant'' are organized as output
of the frame's own postulates rather than as structure imported from
outside --- the organization of these into a single vibrating frame is
not present in the framework's own earlier layers, and not previously
assembled as a single system; the individual elements (elementary
real analysis, linear algebra, the group--velocity splitting, the
zero-bandwidth channel) are textbook material that the frame
organizes rather than merely restates.  To say the claim exactly:
the novelty is claimed at the level of assembly, of mechanism (the
lattice sink-contraction commutator, distinct from the Tomita--Takesaki
modular flow of noncommutativity geometry), and of criterion
structure --- the time count, the width, the dimension bound, the
two branches and their death lines (the structure of criteria the
frame derives from its postulates, as opposed to the values those
criteria take) --- not of new output numbers
(Sec.~IX states honestly: no unconditional new numbers).  The
abstract slogan that time may be read from
noncommutativity has precedent in noncommutative geometry
(Sec.~IX), while the realized criterion --- order-blindness
equivalent to commutativity with the quantitative commutator
$\lambda^2(v_q-v_p)$ --- is new at the level of its quantitative
lattice realization, not of the definitional equivalence itself.  Four claims are established,
each with Lean~4 theorems and zero
\texttt{sorry} (three blocks machine-checked; the empirical
discriminator of the faster-layer branch of the fourth claim is
stated as a proposal, not a machine-checked theorem).  The four
claims regroup and extend the four facts of Sec.~I.
First, \emph{time order is the algebraic shadow
of noncommutativity}: order-blindness is equivalent to commutativity
of the sink operators (AT1--AT2), flat fields are order-blind (AT4),
a gradient makes order observable with commutator
$\lambda^2(v_q-v_p)$ (AT5) --- so time order is \emph{generated by
gradient} (AT10).  The six-round audit of Sec.~IX demotes this
reading: the commutator reports that the two composite orders are
distinguishable, it gives no orientation, no dimension (both sides
are field values) and does not generate time.  Time order survives on massless fields: the
vortex $v=(0,1,1,0)$ is divergence-free, hence massless, yet not flat,
and its gradient keeps the order observable (AT11--AT14).  On the
two-support family --- fields whose fluctuation sits between two
points --- the time observability takes the closed form
$\lambda^2\sqrt{Q_A\,N/(N-1)}$ (N6), proportional to the square root
of the fluctuation energy; the equality is family-specific, not a
general theorem (Sec.~IV).  Second, \emph{a true vibrating coordinate exists iff
two bands have unequal phase velocities}, $\det J=k_1k_2(c_1-c_2)$
(VF2--VF4); the width of the field is the squared beat wavenumber
$(2\pi/\lambda)^2$ (VF8), so ``wide'' and
``fluctuating'' are one parameter; at $k\to0$ the coordinate frame
itself degenerates (VF10) --- the ``instant'' is a simultaneity
surface, not a speed.  Third (demoted by the six-round audit, Sec.~IX),
\emph{the rank bound belongs to the frame}: when all bands share one
$(x,t)$ pair, any three phase functions are linearly dependent
(ND1--ND4) --- a projection-rank statement ($N=3,\dots,1000$ all give
rank~2), not a bound on the dimension of the collective, which is the
number of independent vibration axes (AetherMedium AL4).  Fourth (retracted by the six-round audit, Sec.~IX), the
bandwidth--capacity statement (AB7--AB9) is an internal statement
about a \emph{same-position} time series and does not judge
\emph{cross-position} synchronization; in addition, reach
$\times$ update rate is bounded by $c\varepsilon$ (AB18); a genuinely
faster layer requires an independent cone $c_2=\mathrm{ratio}\cdot c$
whose only empirical discriminator is anisotropy (AB12--AB15).  To
speak of speed one must first have a time order; a time order is
generated by gradient and vanishes in flatness, independently of
masslessness (AT10, AT11--AT14).  The postulates are bold --- space
is a flowing medium, mass is an integer count of strands --- and the
conclusions above are the frame's own output --- as a proposal, with
its boundaries and death lines stated in Sec.~IX --- not a
re-description of known results.
\end{abstract}

\maketitle

% Epigraph provenance: verbatim dialogue supplied by the author,
% spoken by Roy Batty (Rutger Hauer) in Blade Runner (1982), the
% "Tears in Rain" monologue; ellipses retained as supplied. The ZH
% translation lives in the zh mirror file.
\begin{quote}
\itshape
I've\dots\ seen things you people wouldn't believe\dots\\
Attack ships on fire off the shoulder of Orion.\\
I watched c-beams glitter in the dark near the Tannh\"auser Gate.\\
All those\dots\ moments\dots\ will be lost in time, like tears\dots\ in\dots\ rain.\\
Time\dots\ to die\dots
\par
\vspace{0.4em}
--- Roy Batty, \emph{Blade Runner} (1982)
\end{quote}

\noindent\textbf{Note added (2026-10-30).}\emph{The derivation chain of
Secs.~III--VI is kept as written, so that the provenance of the results remains
visible; but six readings have since been judged by a six-round adversarial
audit (Sec.~\ref{sec:audit}) and are retracted or demoted there.  Where the body
below still states them, this note governs: (i) ``the instant implies no
signalling / zero capacity'' (AB7--AB9); (ii) ``time is noncommutativity'';
(iii) ``band number is not dimension'', $\dim=\min(\#\text{bands},d+1)$;
(iv) ``unbounded phase speed implies no faster way''; (v) ``massless $=$
divergence-free''; (vi) ``the aether is $P1$ renamed''.}

\section{Introduction}
\label{sec:intro}

\paragraph{The four facts carried by the intuition.}
Within the flowing-space framework the Yang--Mills structure of
paper~3 \cite{ProjectionPhysics,liu2026yangmills} is proved from
discreteness: mass is an integer count of excited strands,
$m^2(N)=N M_0^2$, so the gap $(0,M_0^2)$ is empty, and the
self-interaction is the noncommutativity of sink contractions,
$(T_pT_q-T_qT_p)v=\lambda^2(v_q-v_p)$ at $p$ (YM3).  If this is the
correct picture, then a natural question follows --- the one this
paper answers: \emph{what is the role of time in a layer of the field
whose fluctuations are so slow that they approach a constant?}  The
question is motivated by the intuition that a ``super-wide'' field ---
a higher collective frequency band --- should act on distant points
faster than light, because the speed of light is the equivalent speed
of space itself.  The intuition carries four distinct facts, which
this paper separates, proves the true ones, and refuses the false
one:

\begin{itemize}
\item \emph{Phase speed is unbounded.}  On a band
$\omega=\omega(k)$, the phase velocity $v_{\mathrm{ph}}=\omega/k$
satisfies $v_{\mathrm{ph}}\to\infty$ as $k\to0$ (AB1--AB2), and the
phase lag between fixed points can be made arbitrarily small (AB3).
``The wider, the faster'' is true at the level of phase.
\item \emph{But instantaneous has zero capacity.}  A zero-bandwidth
mode is a constant sequence (AB7); a message needs bandwidth (AB8) and
the instantaneous channel has capacity exactly $0$ (AB9).  Capacity
grows linearly in dimension but strictly in bandwidth (AB10).
``Instantaneous'' is co-phase --- a simultaneity surface --- not
communication.
\item \emph{Any genuinely faster layer must be an independent
cone.}  Its speed is $c_2=\mathrm{ratio}\cdot c$, a new dimensionless
input, with ratio $=1$ the death line (AB12--AB14), and the only
empirical discriminator between the drag mechanism ($1+v$) and the
hierarchical layer (ratio) is anisotropy (AB15).  The price is a
bound: reach $\times$ update rate $\le c\varepsilon$ (AB18).
\item \emph{Time is the algebraic shadow of the framework's
noncommutativity (YM3).}  Order-blindness is commutativity (AT1--AT2);
flat fields are order-blind (AT4); a gradient makes order observable
with commutator $\lambda^2(v_q-v_p)$ (AT5), so time order is
\emph{generated by gradient} (AT10) and survives on massless fields
(AT11--AT14).  On the two-support family the time observability has
the closed form $\lambda^2\sqrt{Q_A\,N/(N-1)}$ (N6).
\end{itemize}

\paragraph{The claim.}
Within the framework's four postulates --- space is a flowing medium;
mass is an integer count of excited strands; charge is the divergence
of the flow; the four forces are the Leibniz decomposition of flow
momentum --- the following are theorems in Lean~4 with zero
\texttt{sorry} \cite{de2015lean,ProjectionPhysics}: phase speed is
unbounded (AB2); the instantaneous channel has zero capacity (AB9);
reach $\times$ update rate is bounded (AB18); order-blindness is
equivalent to commutativity (AT2); flat implies order-blind (AT4);
gradient makes order observable with commutator $\lambda^2(v_q-v_p)$
(AT5); a flat layer is order-blind (AT4) and a flat massless layer
has no strict sequence (AT9); in the vibrating-coordinate model a
true coordinate exists iff
two bands have unequal phase velocities with
$\det J=k_1k_2(c_1-c_2)$ (VF4), the width is the squared beat wavenumber $(2\pi/\lambda)^2$
(VF8), and the $k\to0$ limit degenerates (VF10);
band number is not dimension --- in $1+1$ dimensions any three bands
are linearly dependent (ND1--ND4).

\paragraph{What this paper is and is not.}
It is a proposal of a new ontology: the universe is a vast collective
of vibration frequencies and dimensions, and physics is the
description process of that field in the vibrating coordinate frame.
The postulates are bold --- space is a flowing medium whose equivalent
speed is $c$; mass is an integer count of excited strands; charge is a
divergence; the four forces are the Leibniz decomposition of flow
momentum --- and the claims are the frame's own output: the
organization of time order from noncommutativity, band number as
dimension, width as fluctuation, the time count as band index times
span, and the instant as a simultaneity surface into a single
vibrating frame is not contained in the program's own earlier papers
and not previously assembled as a single system in the standard
mathematical physics (whose individual elements --- elementary real
analysis, linear
algebra, the group--velocity splitting, the zero-bandwidth channel ---
are textbook material that the frame organizes, not restates).  They
are statements generated by the postulates, the identifications among
them verified by machine checks.  What the paper does not claim: it is not a claim of
faster-than-light signaling, not a claim that the standard framework
is wrong --- it proposes a different substrate, and the hierarchical
layer ($c_2=\mathrm{ratio}\cdot c$) is an open postulate with its own
death line (AB12), not an inference.  Every identification that links
theorems to physics is stated as a proposal of the framework, not
disguised as a theorem.

\paragraph{How to read this paper.}
Sections~II--IV carry the central claim: postulates (Sec.~II); the
layer model of the aether with the two branches and their death lines
(Sec.~III); order as time's algebraic shadow, the new quantitative
bridge, and the correction that masslessness is not flatness
(Sec.~IV).  Section~V introduces the vibrating coordinate frame;
Sec.~VI the band structure: time counting, dispersion and group
velocity, and band number versus dimension.  Section~VII assembles the
chain and relates it to the physical question; Sec.~VIII sets the
chain against the motivating question; Sec.~IX assesses honestly and
states the death lines.

\section{Postulates of the flowing-space framework}
\label{sec:postulates}

The framework is that of paper~1 \cite{ProjectionPhysics}; its
postulates, restated here for the reader, are the following.

\begin{postulate}[Flowing space]
\label{post:flowing}
Space is a flowing medium with a vector equivalent velocity $C$ whose
magnitude is a universal constant $c$: $|C|=c$ at every point.
\end{postulate}

\begin{postulate}[Photons co-move with space]
\label{post:photon}
A photon is entirely carried by the flow: its world line satisfies
$dx=C\,dt$, so its proper time vanishes, $d\tau^2=0$.
\end{postulate}

\begin{postulate}[Three directions]
\label{post:three}
Space carries three independent directions; their Clifford algebra
generates spin.
\end{postulate}

\begin{postulate}[Mass is an anchored count]
\label{post:mass}
Mass is an integer count of excited strands:
$m^2(N)=N\,M_0^2$ with $N$ a nonnegative integer.  A photon has
$N=0$; mass is the anchored resistance of internal motion to space
motion.
\end{postulate}

\paragraph{Reading.}
Nothing in this paper uses a new postulate.  The band picture of
Sec.~III is a describing structure --- phase bands $\{(k,\omega)\}$
--- not a new dynamical postulate; the ``hierarchical layer'' of
branch~II is stated as the only upgrade consistent with the
framework, and its price (the dimensionless ratio) is its own death
line, stated openly as an input (AB12--AB14).

\section{The layer model of the aether: two branches, two death lines}
\label{sec:aether}

\paragraph{The question.}
If the Yang--Mills picture is right, then a synchronizing basis medium
--- an ``aether'' --- should exist that outruns light.  The intuition
is that the universe is a vast collective of vibration frequencies and
dimensions; where local frequencies are low and inconsistent there is
wave phenomena (light speed, spatial ripples), but the super-wide
field is larger, so it can influence distant points ``instantaneously,
far faster than light, because light speed is merely space's equivalent
flow speed.''

\paragraph{The separation.}
The intuition conflates four statements.  The form and the verdict of
each are given in Table~\ref{tab:claims}.

\begin{table}[ht]
\caption{The four statements carried by the ``faster'' intuition, their
forms, and their verdicts.  The boundary between claims~2--3 and
claim~4 is the whole content of Sec.~III: phase speed is not signal
speed, and co-phase is not communication.}
\label{tab:claims}
\begin{ruledtabular}
\begin{tabular}{p{0.12\linewidth}p{0.42\linewidth}p{0.28\linewidth}p{0.08\linewidth}}
\# & Statement & Form & Verdict \\
\hline
1 & Low inconsistent local frequency $\Rightarrow$ waves & $k>0$ and
pointwise distinct $\Rightarrow$ phase gradient (AB6) &
$\checkmark$ ($=$ the definition of a wave)\\
2 & ``Wider'' $\Rightarrow$ larger phase speed, unbounded &
$v_{\mathrm{ph}}=\omega/k$, $k\to0^{+}$ $\Rightarrow$
$v_{\mathrm{ph}}\to\infty$ (AB1, AB2) & $\checkmark$ (elementary
real analysis)\\
3 & The super-wide field ``influences instantaneously'' & fixed $d$:
phase lag arbitrarily small (AB3); $k=0$ mode is everywhere co-phase
(AB5) & $\checkmark$ but co-phase, not signaling\\
4 & Therefore ``far faster than light'' $=$ can transmit faster &
needs bandwidth, needs $k\ne0$, needs $\omega\ne0$ & $\times$\\
\end{tabular}
\end{ruledtabular}
\end{table}

\paragraph{Definitions.}
Let a band be a pair $(k,\omega)$ with dispersion
$\omega=\omega(k)$; the phase is $\theta(x,t)=kx-\omega t$; the phase
velocity is $v_{\mathrm{ph}}(k)=\omega(k)/k$.  Two points are
\emph{synchronized} when their phases agree, $\theta_1=\theta_2$.
A \emph{message} is a time-varying sequence of symbols; a
\emph{channel} has bandwidth $B$ when the frequency support of
sequences it carries is within a band of width $B$.

\begin{theorem}[Phase speed is antitone and unbounded]
\label{thm:ab1}
For a fixed positive $\omega$, $v_{\mathrm{ph}}(k)=\omega/k$ is
strictly decreasing in $k>0$; moreover for every $M$ there exists
$k>0$ with $M<\omega/k$.  Equivalently, \texttt{AB1}: $k_1<k_2
\Rightarrow v_{\mathrm{ph}}(k_2)<v_{\mathrm{ph}}(k_1)$; \texttt{AB2}:
$\forall M\,\exists k>0,\ M<\omega/k$.  Both are proved in Lean with
zero \texttt{sorry}.
\end{theorem}

\begin{theorem}[Co-phase is a well-defined equivalence]
\label{thm:ab4}
The relation ``$\theta_1=\theta_2$ at the same $t$'' is an equivalence
relation (\texttt{AB4}); the constant phase $k=0$ makes every point
co-phase (\texttt{AB5}); a positive gradient makes distinct points
distinct in phase (\texttt{AB6}).
\end{theorem}

\begin{theorem}[The instantaneous channel has zero capacity]
\label{thm:ab9}
A zero-bandwidth sequence is constant (\texttt{AB7}), so a message
requires nonzero bandwidth (\texttt{AB8}), and the capacity of the
$B=0$ channel is exactly zero (\texttt{AB9}).  Capacity is additive in
dimension and strictly increasing in bandwidth, but dimension alone
buys capacity, not single-mode speed (\texttt{AB10}).
\end{theorem}

\begin{theorem}[Reach times update rate is bounded]
\label{thm:ab18}
With phase resolution $\varepsilon>0$, if $k\cdot L<\varepsilon$ then
$c\cdot k<c\cdot\varepsilon/L$ (\texttt{AB18}).  The reach of an
instantaneous layer and its update rate are locked: \emph{the farther
a mode can act, the slower the change it can carry.}
\end{theorem}

\paragraph{The two branches.}
Theorem~\ref{thm:ab9} does not exclude a faster layer; it forces the
faster layer to be a \emph{different} one.

\begin{theorem}[Two branches of ``faster'']
\label{thm:branches}
Branch~I: the super-wide field is the $k\to0$ sector of the same
space field.  It is instantaneous (Theorems~\ref{thm:ab1}--\ref{thm:ab18})
but cannot signal (AB7--AB9); it contributes no new physics beyond
its price, the reach--rate bound
(Theorem~\ref{thm:ab18}).  Branch~II: the super-wide field is an
independent layer with its own cone $c_2=\mathrm{ratio}\cdot c$.  It
is genuinely faster when $\mathrm{ratio}>1$ (\texttt{AB13}); the
degenerate case $\mathrm{ratio}=1$ collapses the two cones and buys no
new speed (\texttt{AB12}), which is the death line; and the two
mechanisms --- coordinate drag $1+v$ and the hierarchical ratio ---
are indistinguishable at a single direction, as witnessed at the
point $v=\tfrac12$ where $\mathrm{dragSpeed}(1/2)=
\mathrm{coarseSpeed}(3/2)$ (\texttt{AB15}); that anisotropy is the
only discriminator is an interpretation of the numerics
(the numerical witness labeled N6-AB in the repository's band
verification), stated as a proposal, not a machine-checked theorem.
\end{theorem}

The logical direction must be stated plainly: the framework proves
``aether $\Rightarrow$ discreteness + noncommutativity $\Rightarrow$
mass gap'', not the reverse.  The hierarchical layer is a reading
consistent with the framework, not an inference from it.

\section{Order as the algebraic shadow of time}
\label{sec:order}

\paragraph{The reading.}
``Before/after'' is the algebraic shadow of time.  It is observable
exactly when two operations do not commute, and the framework has
already proved a noncommutative pair: sink contractions at distinct
points, with exact commutator
\begin{equation}
(T_pT_q-T_qT_p)v \text{ at } p \;=\; \lambda^2(v_q-v_p),
\qquad \text{(\texttt{YM3})}.
\end{equation}
This paper reads YM3 temporally: a nonzero commutator is a scale that
reports which operation came first, so \emph{order is observable};
a zero commutator is a place with no time order.

\begin{theorem}[Order-blindness is commutativity]
\label{thm:orderblind}
A region is \emph{order-blind} when every field gives the same result
under both orders of two contractions; then \texttt{AT1}:
order-blindness is equivalent, by definition, to commutation of the
two contractions on every field; \texttt{AT2} (the machine-checked
form): the region is not order-blind iff there exists a field on which
the two orders give different fields.
\end{theorem}

\begin{theorem}[Flat fields are order-blind]
\label{thm:flat}
A constant field is left unchanged by any single-point contraction
(\texttt{AT3}), so on a flat region every two contractions commute and
the region is order-blind (\texttt{AT4}).  Flatness, not
masslessness, is the driver: masslessness does not imply flatness
(Theorem~\ref{thm:contra}).
\end{theorem}

\begin{theorem}[Gradient makes order observable]
\label{thm:gradient}
If $v_q\ne v_p$ and $\lambda\ne0$, the commutator at $p$ is nonzero
(\texttt{AT5}, reusing YM3); non-flatness guarantees such a pair
(\texttt{AT6}).  Hence \texttt{AT10}: \emph{a gradient makes "which
came first" resolvable --- time order is generated by gradient.}
\end{theorem}

\begin{theorem}[A flat massless layer has no strict sequence]
\label{thm:massless}
A co-moving proper interval vanishes: $dx=c\,dt \Rightarrow
d\tau^2=0$ (\texttt{AT7}, restating SM1).  When every interval in a
layer vanishes, there is no strict order of any two events
(\texttt{AT8}); combining with Theorem~\ref{thm:flat}: a flat
massless layer is order-blind and has no strict time order
(\texttt{AT9}; here ``massless'' reads as ``all intervals in the
layer vanish,'' a reading of the frame --- the Lean statement's
hypothesis is the vanishing itself).  Masslessness alone does not
provide this: the
vortex of Theorem~\ref{thm:contra} is massless, non-flat, and keeps
its order observable (AT13).
\end{theorem}

\begin{theorem}[The quantitative bridge]
\label{thm:bridge}
For any configuration the commutator field is constant with value
$\lambda^2(v_q-v_p)$ --- directly from the sink-contraction
arithmetic (the Lean statement YM3 fixes the value at $p$; the
constancy over the region is the same two-point arithmetic, not a
separately formalized Lean theorem) --- so the time observability
defined as the root mean squared commutator over its region is
$|\lambda^2(v_q-v_p)|$.  On the two-support family --- fields whose
only nonzero fluctuation sits between two points (including $N=2$) ---
this equals
\begin{equation}
\text{time observability }
\;=\; \lambda^2\,\sqrt{\frac{Q_A\,N}{N-1}},
\label{eq:timeobs}
\end{equation}
with $Q_A$ the fluctuation energy of the region (the framework's
gravity primitive: $Q_A=0$ iff the region is flat) and $N$ the number
of sites; the identity is verified to machine precision on 15
deterministic configurations of this family (max relative discrepancy
$2.17\times10^{-16}$, \texttt{N6}).  The equality is
family-specific, not a general theorem.  The reference magnitudes
(deviations relative to the bridge value $B$ and to the measured
value $M$ respectively) differ in the two counterexamples: on the
linear ramp $v=(0,1,2)$ the bridge \emph{underestimates} the
measurement ($M/B=1.155$, i.e. the measurement exceeds the bridge by
$15.5\%$ of $B$, equivalently $13.4\%$ of $M$); on the vortex of
Theorem~\ref{thm:contra} (with $p,q$ at unequal entries) the bridge
\emph{overestimates} the measurement ($B/M=1.155$, i.e. the bridge
exceeds the measurement by $15.5\%$ of $M$, equivalently $13.4\%$
of $B$).  The two deviations mirror each other because the two
configurations have reciprocal ratio structures
$\sqrt{Q_A N/(N-1)}/|v_q-v_p| = \sqrt3/2$ (ramp) and
$\sqrt{4/3}$ (vortex), whose product is $1$.  What holds generally
is the qualitative content:
when the fluctuation vanishes, $Q_A=0$, the field is flat, and the
time order disappears exactly (\texttt{N1}); the quantitative
$\sqrt{Q_A}$ proportionality is exact only within the two-support
family.
\end{theorem}

\paragraph{The correction: masslessness is not flatness.}
The first reading of this section conflated two attributes.  The
user's correction is a type-level distinction: \emph{masslessness} is a
property of the \emph{source} --- it produces no distortion of the
surrounding space (no divergence) --- whereas \emph{flatness} is a
property of the \emph{field} (constant value, $Q_A=0$).  The two are
not equivalent, and the massless layer does not require flatness:

\begin{theorem}[Masslessness is not flatness]
\label{thm:contra}
A constant field has zero discrete divergence on any four-point loop
(\texttt{AT11}: flat implies distortion-free).  The vortex field
$v=(0,1,1,0)$ has zero divergence (\texttt{AT12}: massless) yet is
not flat, with curl $=2$; on it, the time order is still observable
since $v_0\ne v_1$ gives a nonzero commutator (\texttt{AT13}).
Therefore flatness is a strict subset of masslessness
(\texttt{AT14}): flat is a local special case, and the reading
``massless $\Rightarrow$ no time order'' is withdrawn.
\end{theorem}

What survives: the instant is not ``faster than light'', but the
reason is the flat/no-gradient branch (Theorems~\ref{thm:flat}--\ref{thm:gradient}),
not masslessness.  The interface with Sec.~III is unchanged:
zero bandwidth is zero capacity.

\section{The vibrating coordinate frame}
\label{sec:vibration}

The zero-capacity line closes the channel reading of the instant; we
now build the coordinate frame in which the field itself is
described.  \paragraph{The model.}
A field is described by a family of bands $\{(k,\omega,w)\}$, each
with phase $\theta(x,t)=kx-\omega t$ and positive weight $w$, and
linear dispersion $\omega=c\cdot k$.  The description process is three
steps: (i) field $\to$ vibrating coordinates, (ii) ``width'' as the
variance of the wave-vector distribution, (iii) the super-wide limit
$k\to0$.

\begin{theorem}[A single mode carries one dimension]
\label{thm:vf1}
A single mode with linear dispersion depends only on the combination
$\xi=x-ct$: $\theta(x+cs,t+s)=\theta(x,t)$ (\texttt{VF1}).  One mode
alone has no coordinate structure; ``width'' is not definable within
it.
\end{theorem}

\begin{theorem}[True coordinates need two unequal speeds]
\label{thm:vf2}
The Jacobian of the map $(x,t)\mapsto(\theta_1,\theta_2)$ for two
bands is a closed form
\begin{equation}
\det J \;=\; k_1k_2\,(c_1-c_2),
\label{eq:jac}
\end{equation}
(\texttt{VF2}).  Two bands with the same phase speed have
$\det J=0$ (\texttt{VF3}); two bands with unequal phase speeds and
nonzero wave numbers have $\det J\ne0$ (\texttt{VF4}).  \emph{A true
vibrating coordinate exists if and only if the phase speeds differ.}
\end{theorem}

\begin{theorem}[Width is the squared beat wavenumber]\label{thm:vf8}
The width of the field, the variance of the wave-vector distribution,
has the closed form
\begin{equation}
\mathrm{var}
\;=\;\frac{w_1w_2}{(w_1+w_2)^2}\,(k_1-k_2)^2
\;=\;C\,\left(\frac{2\pi}{\lambda}\right)^2,
\label{eq:var}
\end{equation}
with $\lambda=2\pi/|k_1-k_2|$ the beat wavelength and
$C=w_1w_2/(w_1+w_2)^2$ (\texttt{VF5}, \texttt{VF7}, \texttt{VF8}).
Width is zero iff the wave numbers coincide (\texttt{VF6}).
\emph{``Wide'' and ``fluctuating'' are two names for one parameter.}
\end{theorem}

\begin{theorem}[The super-wide limit kills the coordinate]
\label{thm:vf10}
At $k_1=k_2=0$, every phase vanishes, $\theta\equiv0$
(\texttt{VF9}, \texttt{VF10}); the variance is zero, and the Jacobian
is zero.  The vibrating coordinate system itself degenerates: in the
super-wide limit there are no coordinates, hence no ``faster/slower''
comparison, hence no ``instant'' to speak of.
\end{theorem}

\paragraph{Interpretation.}
``Wide'' is a direction in $\Delta k$; ``fast'' is a different
direction in $\Delta c$; the illusion of superluminality arises from
merging them into one axis.  In the super-wide limit ($\Delta k\to0$)
both axes collapse together --- the coordinates fail, and the instant
is not a speed but the absence of speed.

\section{Band structure: time counting, dispersion, and dimension}
\label{sec:bands}

\paragraph{The time-count interface.}
The framework's closure counting (VBC: a time unit is one closure of
the phase) and the band picture (a wave number measures a wavelength)
are the same number written twice: one closure of the phase is a
displacement of one wavelength, so

\begin{theorem}[Time count is band index times span over $2\pi$]
\label{thm:tc}
On a span $L$ the number of closures is
\begin{equation}
n(k,L)=\frac{kL}{2\pi},
\qquad
n=\rho\,L\quad\text{with}\quad \rho=\frac{k}{2\pi}=\frac1\lambda,
\label{eq:ntc}
\end{equation}
(\texttt{TC1}).  The count is additive in the span (\texttt{TC2}); one
wavelength gives exactly one closure, $n(k,2\pi/k)=1$ (\texttt{TC3}).
The count depends on the chosen band: $k_1\ne k_2$ admits a span where
the counts differ (\texttt{TC4}) --- \emph{there is no absolute time
count, only a count on a band.}  The global mode $k=0$ counts zero
(\texttt{TC5}), and the count-difference density is the inverse beat
wavelength,
$|\Delta k|/2\pi=1/\lambda_{\mathrm{beat}}$ (\texttt{TC6}): the time
count and the ``width'' of Sec.~V are the same scale.
\end{theorem}

\paragraph{General dispersion: the criterion survives, the speed
splits.}
The vibrating-coordinate criterion of Sec.~V was
$\det J=k_1k_2(c_1-c_2)$, special to linear dispersion.  Under a
general dispersion the criterion is

\begin{theorem}[Degeneracy iff equal phase speeds]
\label{thm:gd}
The generalized Jacobian determinant is
\begin{equation}
\det\nolimits_2(m_1,m_2)
\;=\; k_1\omega_2-k_2\omega_1
\;=\; k_1k_2\,(v_{{\mathrm{ph}},2}-v_{{\mathrm{ph}},1}),
\label{eq:det2}
\end{equation}
(\texttt{GD1}); for nonzero wave numbers it vanishes iff the two
bands have equal phase velocities (\texttt{GD2}; the $k=0$ degeneracy
is covered by VF10).  Under the quadratic family
$\omega=ck+\beta k^2$, the determinant is $k_1k_2\beta(k_2-k_1)$
(a direct corollary of GD1, verified numerically as N9), with
$v_{\mathrm{ph}}=c+\beta k$ and group velocity
$v_g=d\omega/dk=c+2\beta k$.  In linear dispersion
$v_g=v_{\mathrm{ph}}$ (\texttt{GD7}); under nonlinearity
$v_g\ne v_{\mathrm{ph}}$ (\texttt{GD8}); the group velocity gap is
\begin{equation}
|\Delta v_g|\,\lambda_{\mathrm{beat}}=4\pi|\beta|
\label{eq:gvgap}
\end{equation}
(\texttt{GD9}, absolute form GD9b): \emph{the wider the field, the slower the ripple
(envelope); as $\lambda\to\infty$ the envelope freezes.}  As in
Sec.~IX, the $k$-form $|\Delta v_g|=2|\beta||\Delta k|$ is the
machine-checked statement (GD9/GD9b in Lean); the $\lambda$-form
displayed is the same identity rewritten with
$\lambda_{\mathrm{beat}}=2\pi/|\Delta k|$, verified numerically
(N12).  The same
statement as Theorem~\ref{thm:vf10} in group-velocity language.
\end{theorem}

\paragraph{Band number is not dimension.}
In $1+1$ dimensions, $(k,\omega)\in\mathbb{R}^2$, so any three bands
are linearly dependent.

\begin{theorem}[Three bands are redundant in 1+1 dimensions]
\label{thm:nd}
If the first two bands are non-degenerate,
$\det\nolimits_2(m_1,m_2)\ne0$, there exist coefficients $(a,b)$
such that a third band satisfies $k_3=ak_1+bk_2$ and
$\omega_3=a\omega_1+b\omega_2$ (\texttt{ND1}, Cramer's rule); the third phase is then the same combination for all
$(x,t)$, $\theta_3=a\theta_1+b\theta_2$ (\texttt{ND2}); and even its
time count is redundant, $n_3(L)=an_1(L)+bn_2(L)$ (\texttt{ND4}),
with the rank of the three bands being $2$ (\texttt{ND3}).  \emph{Adding bands does
not add dimensions} is the theorem (rank~$\le2$ in $1+1$
dimensions); the rank bound itself is a pure linear-algebra theorem
in any dimension (rank $\le d+1$ on $d+1$ linear forms), while the
formula $\dim = \min(\#\text{bands},\,d+1)$ that reads the bound as
\emph{the dimension of the vibrating coordinate} is a reading of the
framework, stated
as a proposal (formalized only for $1+1$).  To obtain a third
independent coordinate one must raise the spatial dimension, not add
a collinear band.
\end{theorem}

\section{The assembled chain}
\label{sec:chain}

The four facts of the introduction are not one theorem, but they
share one logical skeleton.  The material is organized in three
disjoint machine-checked clusters --- {AB}, {AT+YM+SM},
{VF+TC+GD+ND} --- and one \emph{dictionary seam}: an explicit
correspondence among three proved criteria (VF4, AT5 and GD2)
that lets the two participating clusters speak the same language.  The {AB} cluster does
not participate in this seam; it is linked to the rest at the level
of commentary, not of theorem.  The seam of this section is the identity of the non-degeneracy
criterion in the VF, AT and GD vocabularies.

\paragraph{The dictionary seam (explicit correspondence of proved
criteria).}
The vibrating frame proves (VF4, machine-checked): two bands with
nonzero wave numbers give a true coordinate iff their phase speeds
differ, $k_1k_2\ne0 \wedge (\det J=k_1k_2(c_1-c_2)\ne0
\Leftrightarrow c_1\ne c_2)$.  The order theorem proves (AT5,
machine-checked, reusing YM3): a gradient makes the order observable,
with commutator $\lambda\ne0 \wedge (\lambda^2(v_q-v_p)\ne0
\Leftrightarrow v_q\ne v_p)$.  The general dispersion theorem (GD1,
machine-checked) writes the same criterion as
$\det_2=k_1k_2(v_{{\mathrm{ph}},2}-v_{{\mathrm{ph}},1})$.  The
seam is the shared shape of these three proved statements --- now
itself machine-checked (Seam module, Lean): the lemma
${\tt seam\_iff}: a\ne0 \rightarrow (a(x-y)\ne0 \leftrightarrow
x\ne y)$ is proved once, and VF4, AT5 and GD2 are declared as its
three instances.  The seam is not a new theorem about a new object:
it is the same criterion written in three vocabularies, and the
identification that ties the vocabularies --- the \emph{identification of the field value at a point with a band phase
speed} ($v_i:=c_i$; an identification of the framework, stated as a
proposal, not a theorem) --- remains a proposal even though the
shared shape is machine-checked.  The three statements of the same
shape, with the prerequisites $k_1k_2\ne0$ and $\lambda\ne0$ made
explicit:
\begin{equation}
\begin{aligned}
\text{coordinate non-degenerate} &\;\Leftrightarrow\; c_1\ne c_2
\;\Leftrightarrow\; \text{order observable} \\
&\;\Leftrightarrow\; v_q\ne v_p \;\Leftrightarrow\;
\det J\ne0 \;\Leftrightarrow\; \text{commutator at } p \ne0 ,
\qquad (k_1k_2\ne0,\; \lambda\ne0)
\end{aligned}
\label{eq:seam}
\end{equation}
The common shape is the point: each of the three criteria
reduces to ``the coefficient is nonzero and the velocity difference
is nonzero.''  This is what the ``instant'' intuition is doing when
it fuses ``wider'' with ``faster'': it collapses the two axes
$\Delta k$ and $\Delta c$ of the vibrating frame into one.  The
seam makes the fusion explicit: the shared shape is machine-checked
(Seam module), and the identification $v_i:=c_i$ that ties the
vocabularies remains a proposal of the framework.

The instant is a simultaneity surface, not a speed.  A speed is a
rate over a time order; a time order is generated by a gradient
(Theorem~\ref{thm:gradient}), and the gradient is exactly what a flat
field lacks (Theorem~\ref{thm:flat}): the wider the field, the
flatter the phase, the smaller the time order, and in the $k\to0$
limit the vibrating coordinates collapse
(Theorem~\ref{thm:vf10}) and the instantaneous channel has zero
capacity (Theorem~\ref{thm:ab9}).  The only upgrade to a genuine
speed is a new layer with its own cone and its own dimensionless ratio
(AB12--AB15), whose empirical discriminator is anisotropy.

\section{Relation to the question as posed}
\label{sec:relation}

This paper does not prove faster-than-light signaling, and does not
claim the universe \emph{is} a vibrating collective in the sense of a
physical aether; the ontology of the title is claimed at the level of
the description process --- physics as the description of the field in
the vibrating coordinate frame --- not as an assertion about an
underlying substance.  Its relation to the motivating question is the
following.

\begin{enumerate}
\item \emph{How the statement differs from the intuition.}
The intuition is a single claim: ``wider, hence faster.''  This paper
splits it into four statements (Table~\ref{tab:claims}), three true at
the level of phase and one false at the level of signal.  The
boundary is the mathematical content: a message is a time-varying
sequence, and a time-varying sequence requires nonzero bandwidth; the
constant mode is a simultaneity surface.
\item \emph{Can we derive the intuition?}
Not as stated.  The ``genuinely faster'' branch requires a new
dimensionless ratio, and the framework has no equation forcing it;
that is the second input gap, of the same kind as $\lambda$ and
$M_0$.  The death line is written: if the ratio is derivable from the
postulates, branch~II stops being a new postulate and the gap closes.
\item \emph{Why can the framework say something here at all?}
Because the framework has a native noncommutativity (YM3) and a
native count (VBC closures), and both are theorems, not postulates.
The time order is the reading of a proved commutator; the band count
is the reading of a proved closure count.  The paper does not add a
time postulate.
\item \emph{Why the intuition is dangerous as stated.}
Because in a general dispersion the phase speed is not the signal
speed (GD8), and reading ``unbounded phase speed'' as ``faster
signal'' is exactly the sleight that superluminal claims are built on;
AB7--AB9 are the formal refusal.
\end{enumerate}

\section{Six-round adversarial audit (2026-10-30)}
\label{sec:audit}

Sections~III--VI were written first and then audited by six independent
adversarial lines under a frozen contract whose criterion was
\emph{fidelity to the framework's own derivation order}, with standard-physics
no-go statements (no-signalling, light cones, Shannon capacity, QFT axioms)
explicitly barred as judging criteria.  The verdicts were: two claims
\emph{retracted}, four \emph{demoted}, none left standing unchanged.

\begin{enumerate}
\item \textbf{(R1, retracted) ``instant implies no signalling''.}
AB7--AB8 use $s:\mathbb{Z}\to\alpha$, a time series at \emph{one} position.
The framework's claim is at \emph{one instant, two positions}.  The axes were
swapped.  The replacement object is the medium fibre
$\mathrm{MediumFiber}\,L := \lambda x\,y.\, L x = L y$ (an equivalence
relation, not an order and not a channel), and the decisive lemma is that
spatial co-phase does \emph{not} imply temporal constancy
(\texttt{spatial\_const\_not\_time\_const}).
\item \textbf{(R2, demoted) ``time is noncommutativity''.}
The sink commutator $\lambda^{2}(v_q-v_p)$ reports that two composite orders
are distinguishable; it gives no orientation (it is symmetric in $p,q$), no
dimension (both sides are field values), and there is no map from field values
to seconds.  Time's interface is the band index and the closure count
$n=kL/2\pi$ (TC1--TC6, VBC4--VBC5).
\item \textbf{(R3, demoted) ``band number is not dimension''.}
The bound $\mathrm{rank}\le\min(\#\text{bands},d+1)$ holds only when all bands
share one $(x,t)$ pair, and it is insensitive to the band count
($N=3,\dots,1000$ give rank~2).  It is a projection-rank (factorization)
statement whose positive content is \emph{phase locking}: for $N>2$ shared-frame
bands there is a non-trivial linear relation.  Dimension is the number of
independent vibration axes, an input.
\item \textbf{(R4, demoted) ``group velocity corrects the phase-speed reading''.}
$c$ is a medium-layer constant (\texttt{MediumFlow.c}, no $k$); $v_{ph}$ and
$v_g$ are wave-layer descriptions of the same medium.  GD8 is a true statement
inside the wave layer and does not judge the medium-layer claim.
\item \textbf{(R5, demoted) ``massless $=$ divergence-free''.}
The divergence is the \emph{charge}-layer quantity (P3: charge $=$ divergence);
mass is the strand count / anchoring (P2, MC1: $m^{2}=N M_0^{2}$).  Moreover the
vortex $v=(0,1,1,0)$ has a gradient, hence is \emph{massive} by MC1: it witnesses
``sourceless but not flat'', not ``massless but not flat''.
\item \textbf{(R6, retracted) ``the aether is just $P1$ renamed''.}
With $\mathcal{A}:=C$ a definitional identity, the statement ``$\mathcal{A}$
exceeds $C$'' is not expressible and every predicate is blind to the difference
(Leibniz).  The medium must be a \emph{fibre} of the field map:
$\mathrm{observe}:\text{Medium}\to\text{Field}$ is not injective, and no
functional of the field alone reads the medium's synchronization speed scale
(\texttt{no\_field\_functional\_reads\_speedScale}).
\end{enumerate}

What survives is the framework's own positive content, now stated at the correct
layer (\texttt{Explorations/AetherMedium.lean}, AL1--AL7, zero \texttt{sorry}):
the field cannot fix the medium (AL1, AL7); synchronization is an equivalence
relation on the medium layer (AL2); the time coordinate is supplied by the band,
so no band means \emph{no period}, not ``period zero'' (AL3); the rank bound
belongs to the frame (AL4); masslessness is absence of anchoring and is
independent of the field's shape (AL5); no single band carries the medium speed
(AL6).  The dynamics of the wide field --- the positive task of the framework ---
is \emph{not} supplied here and remains open.

\section{Honest assessment}
\label{sec:honest}

\begin{table}[h]
\centering
\caption{The four facts of Sec.~I, the four blocks of claims that
regroup and extend them, the verification status of each, and the
death line attached to it.}
\begin{ruledtabular}
\begin{tabular}{llll}
\multicolumn{1}{c}{Fact (Sec.~I)} & \multicolumn{1}{c}{Block} &
\multicolumn{1}{c}{Status} & \multicolumn{1}{c}{Death line}\\
\hline
phase speed unbounded (AB1--AB2) & the instant, its capacity, and
the faster layer (AB1--AB18) & machine-checked, zero sorry &
a collective mode at $k=0$ carries a time-varying symbol
(branch~I)\\
instantaneous channel, zero capacity (AB7--AB9) & same block &
machine-checked & same death line, branch~I\\
faster layer needs independent cone (AB12--AB18) & same block &
discriminator (AB15) is a proposal, not a theorem &
ratio derivable from the postulates, or a fast factor neither drag
nor ratio (branch~II)\\
time is the shadow of noncommutativity (YM3) & time order = algebraic
shadow of noncommutativity (AT1--AT14) & machine-checked &
two sink operators fail to commute on a flat field\\
& true coordinate iff phase speeds differ (VF2--VF10) & machine-checked &
phase speeds equal for nonzero wave numbers\\
& band number is not dimension (ND1--ND4) & machine-checked &
phase rank fails to match physical dimension under $v_i:=c_i$\\
\end{tabular}
\end{ruledtabular}
\end{table}

\begin{itemize}
\item \emph{Mathematical content.}  All theorems of Secs.~III--VI are
elementary real analysis, linear algebra, and order relations
(machine-checked in Lean~4, zero \texttt{sorry}); the boundaries are
the ones stated: \texttt{GD5} (the group velocity is the derivative)
is not Lean-proved (a formalization-engineering gap in mathlib
v4.35), with a numerical substitute; the dimension cap is covered in
$1+1$ dimensions only.
\item \emph{Structural correspondence.}  The identification of the
time order with noncommutativity is a reading of the framework, not a
theorem that ``time equals noncommutativity''; the idea that time
can be read off from noncommutativity has precedents in
noncommutative geometry (Connes~\cite{connes1994ncg}, whose modular
automorphism is a very different mechanism --- the Tomita--Takesaki
flow on von Neumann algebras, against our lattice sink-contraction
commutator; we cite it as precedent of the idea, not of the
mechanism), which we
acknowledge while the lattice realization here is the frame's own;
the vibrating frame is
a two-band linear-dispersion model choice; the hierarchical layer is
a reading consistent with, but not implied by, the framework.
\item \emph{Numerical matching.}  All numerical assertions in
Secs.~IV--VI are machine-precision checks of identities
(max residual $\sim 10^{-15}$); the time-observability formula
(Eq.~\ref{eq:timeobs}) is an algebraic combination on a restricted
family, not an unconditional prediction (its conditional status is
stated in the sleeping-discriminator entry below).
\item \emph{Predictions, stated honestly in three layers.}
\emph{(i) Unconditional predictions: none.}  This paper proves
boundaries and gives criteria shapes, not new numbers; beyond the
organizing identities, the only new quantitative closed form is
Eq.~\ref{eq:timeobs}, which is an
algebraic combination on a restricted family, not a prediction.
\emph{(ii) Conditional discriminative predictions: three.}  If a
genuinely faster-than-$c$ factor exists, the two branches of
Sec.~III carry distinct angular signatures with zero new input ---
coordinate drag is anisotropic ($1+v(\theta)$) while the hierarchical
ratio is isotropic --- so an observed angular pattern that matches
neither kills both branches (AB15).  If a collective mode at
$k=0$ ever carries a time-varying resolvable symbol, branch~I dies
(AB7--AB9).  And in any medium with quadratic dispersion
$\omega=ck+\beta k^2$, the group-velocity gap and the beat
wavelength are locked by $|\Delta v_g|\,\lambda_{\mathrm{beat}}
=4\pi|\beta|$ (GD9: the $k$-form $|\Delta v_g|=2|\beta||\Delta k|$
is machine-checked; the $\lambda$-form is the same identity rewritten
with $\lambda_{\mathrm{beat}}=2\pi/|\Delta k|$, verified numerically
as N12): a measured pair
$(\Delta v_g,\lambda_{\mathrm{beat}})$ determines $\beta$, and any
deviation is a framework-internal falsifier.  \emph{(iii) Extension-layer predictions (with new
inputs): two.}  The framework's own experimental program already
lists the signatures of its un-derived postulates --- the cyclotron
anomaly $R_{ci}=1/(1-\mu)-1$ (zero-hypothesis $R_{ci}=0$) and the
spin-polarized splitting
$\Delta R_{ci}=-2\varepsilon/((1-\mu)(1-\varepsilon^2))$ --- which
are conditional on the inputs $\mu,\varepsilon$ (SA8--SA9); we
record them here as the program's empirical contact, not as claims of
this paper.
\item \emph{The second input gap is unchanged.}  The ratio
(AB14), $\lambda$, and $M_0$ remain inputs; no equation in the
framework derives them.  Eq.~\ref{eq:timeobs} is here registered as a
\emph{sleeping discriminator}: it is an exact structural identity on
the two-support family (not a prediction) that becomes a numeric
discriminator the moment $Q_A$ acquires an observational dictionary;
the converse mapping --- a measurement procedure for the
fluctuation energy $Q_A$ --- is a named open problem of the
framework (first step: anchor the dimension of $Q_A$ with the
Yang--Mills scale $\lambda^2 M_0^2$).
\item \emph{Death lines.}  Branch~I dies if a collective mode at
$k=0$ carries a time-varying resolvable symbol; branch~II dies if the
ratio is derivable from the postulates or if an observed ``fast
factor'' is neither drag nor ratio; the flatness reading dies if two
sink operators fail to commute on a flat field; the dimension cap
(the reading $\dim=\min(\#\text{bands},\,d+1)$) dies if, in any
concrete model, the phase rank fails to match the physical dimension
under the identification $v_i:=c_i$ --- the rank bound itself is a
theorem in any dimension (rank $\le d+1$ in $d+1$ dimensions), so
the death line sits in the identification layer, which a concrete
model can falsify, not in the rank theorem.
\item \emph{Aether, synchronization, and signaling (firewall).}  Within the
framework the ``aether'' (P1) permits a simultaneity surface
(co-phase, $d\tau=0$) but, by AB7--AB9, carries zero information
capacity; it is branch~I, not a channel.  A genuinely faster layer
(branch~II, $c_2=\mathrm{ratio}\cdot c$) would require a new dimensionless
input and is stated openly as an un-derived postulate, not an inference (AB12
is the death line).  The program's claim is therefore exactly the boundary
proved here: the intuitive ``faster sync'' is either co-phase (free, no
signal) or a new postulate (costs $\mathrm{ratio}$); we assert neither faster
signaling nor its denial as physics.  No statement in this paper upgrades the
simultaneity surface to a communication channel.
\end{itemize}

\section{Conclusion}
\label{sec:conclusion}

Within the flowing-space framework, the ``instantaneous, faster than
light'' intuition is separated into four statements, three of which
are proved true at the level of phase and one of which is refused at
the level of signal, each of the claims carried by the paper
established with machine-checked theorems.  The phase velocity is
unbounded; the
instantaneous channel has zero capacity; the only genuine speed is a
new layer with a new dimensionless ratio and anisotropy as its only
discriminator; and time itself is the algebraic shadow of the
framework's noncommutativity, generated by gradient, absent in
flatness, and on the two-support family it scales with the square
root of the fluctuation.
The wider the field, the flatter the phase; the flatter the phase, the
smaller the time order; and in the super-wide limit there is neither
coordinate nor channel --- only a simultaneity surface.  The instant
is not a speed.  The ontology proposed by this paper is the vibrating
coordinate frame itself: time order, band number, dimension, width,
and the instant are outputs of its postulates, machine-checked in
Lean, and the substrate it proposes --- boldly postulated, honestly
bounded --- is an assembly the program's own earlier papers do not
contain and the standard mathematical physics does not offer as a
single system (its elements being textbook material, its
time-from-noncommutativity idea having precedents in noncommutative
geometry).

\section*{Conflict of interest}
The author declares no competing interests.

\section*{Data availability}
All theorems are machine-checked in Lean~4 (zero \texttt{sorry});
all numerical claims are generated by the repository's verification
scripts with reproducible random seeds.  Source code and the
full theorem statements are available in the ProjectionPhysics
repository.

\bibliography{projection-physics}

\end{document}